\subsection{SOME ELEMENTS OF SL(2, k)}

For any \(t\in k^{*}\) , put

\[
\begin{array}{l} \theta (t) = e ^ {t X _ {+}} e ^ {t ^ {- 1} X _ {-}} e ^ {t X _ {+}} \\ \qquad = \left( \begin{array}{c c} 1 & t \\ 0 & 1 \end{array} \right) \left( \begin{array}{c c} 1 & 0 \\ - t ^ {- 1} & 1 \end{array} \right) \left( \begin{array}{c c} 1 & t \\ 0 & 1 \end{array} \right) \\ \qquad = \left( \begin{array}{c c} 0 & t \\ - t ^ {- 1} & 0 \end{array} \right) \\ \qquad = e ^ {t ^ {- 1} X _ {-}} e ^ {t X _ {+}} e ^ {t ^ {- 1} X _ {-}}. \end{array}
\]

With the notations of no. 3,

\[
\theta (t) u = - t ^ {- 1} v \quad \theta (t) v = t u
\]

so

\[
\theta (t) e _ {n} ^ {(m)} = (- 1) ^ {m - n} t ^ {2 n - m} e _ {m - n} ^ {(m)}.\tag{4}
\]

Hence, the element \(\theta(t)^2 = \left( \begin{array}{cc} -1 & 0 \\ 0 & -1 \end{array} \right)\) operates by \((-1)^m\) on \(\mathrm{V}(m)\) . If \(\mathrm{E}\) is an odd-dimensional simple \(\mathfrak{sl}(2,k)\) -module, \(\theta(t)_{\mathrm{E}}\) is thus an involutive automorphism of the vector space \(\mathrm{E}\) . In particular, taking \(\mathrm{E}\) to be the adjoint representation:

\[
\theta (t) _ {\mathrm{E}} X _ {+} = t ^ {- 2} X _ {-} \quad \theta (t) _ {\mathrm{E}} X _ {-} = t ^ {2} X _ {+} \quad \theta (t) _ {\mathrm{E}} H = - H\tag{5}
\]

so that \(\theta(1)_{\mathrm{E}} = \theta(-1)_{\mathrm{E}}\) is the canonical involution of \(\mathfrak{sl}(2,k)\) .

For any \(t\in k^{*}\) , put

\[
h (t) = \left( \begin{array}{c c} t & 0 \\ 0 & t ^ {- 1} \end{array} \right) = \theta (t) \theta (- 1).
\]

Then \(h(t)u = tu\) , \(h(t)v = t^{-1}v\) , so

\[
h (t) e _ {n} ^ {(m)} = t ^ {m - 2 n} e _ {n} ^ {(m)}.\tag{6}
\]

PROPOSITION 6. Let E be a finite dimensional \(\mathfrak{sl}(2,k)\) -module, and \(t\in k^{*}\) . Let \(\mathrm{E}_p\) be the set of elements of E of weight \(p\) .

\begin{enumerate}
\def\labelenumi{(\roman{enumi})}
\item
  \(\theta(t)_{\mathrm{E}}|\mathrm{E}_{p}\) is a bijection from \(\mathrm{E}_p\) to \(\mathrm{E}_{-p}\) .
\item
  \(h(t)_{\mathrm{E}}|\mathrm{E}_p\) is the homothety with ratio \(t^p\) on \(\mathrm{E}_p\) .
\end{enumerate}

If \(\mathrm{E} = \mathrm{V}(n)\) , the proposition follows from formulas (4) and (6). The general case follows from Th. 1.

COROLLARY. Let \(E = E' \oplus E''\) be the decomposition of E defined in the Cor. of Prop. 2. The element \(\begin{pmatrix}-1 & 0 \\ 0 & -1\end{pmatrix}\) of \(\mathbf{SL}(2, k)\) operates by +1 on \(E'\) and by -1 on \(E''\) .

This follows from (ii), applied to \(t = -1\) .

